%% ma: L10-B13-0097
%% lop: 10
%% chuong: 5
%% bai: 13
%% dang: 10.13.13
%% loai: DS
%% muc: [TH, TH, NB, TH]
%% dap_an: "ĐSSS"
%% nguon: "chuyen-de-cac-so-dac-trung-cua-mau-so-lieu-khong-ghep-nhom-toan-10-chuong-trinh-moi.pdf (GV. Phan Nhật Linh) (P03), Bài 02 Dạng 2 Phần II Câu 1"
%% kien_thuc: "Lớp 10 Bài 13 (số trung bình, trung vị, tứ phân vị, mốt)"
%% trang_thai: chinh-thuc
%% ghi_chu: "Tự duyệt theo yêu cầu giáo viên 22/10/2026."
\begin{ex}
Một cửa hàng bún đậu vừa khai trương thống kê lượng khách tới quán trong $7$ ngày đầu như sau:
\begin{center}
\begin{tabular}{|c|c|c|c|c|c|c|c|}\hline
Ngày & $1$ & $2$ & $3$ & $4$ & $5$ & $6$ & $7$\\\hline
Số khách & $575$ & $454$ & $400$ & $325$ & $351$ & $333$ & $412$\\\hline
\end{tabular}
\end{center}
Xét tính đúng sai của các khẳng định sau:
\choiceTF{\True Số trung bình của mẫu số liệu, làm tròn đến hàng phần trăm, là $407{,}14$}{Trung vị của mẫu số liệu là $263$}{Ngày 1 là mốt của mẫu số liệu}{Nếu ngày 6 có $400$ lượt khách thì mốt của mẫu số liệu là ngày 3}
\loigiai{a) Đúng: $\overline x=\dfrac{2850}7\approx407{,}14$.

b) Sai: sắp xếp $325$; $333$; $351$; $400$; $412$; $454$; $575$, $M_e=400$.

c) Sai: mốt là một giá trị của mẫu (số khách), không phải ngày; mẫu này mọi giá trị đều xuất hiện một lần.

d) Sai: khi đó mốt là $400$ (lượt khách), không phải ``ngày 3''.}
\end{ex}
